\documentclass[11pt]{article}
\usepackage[margin=1in]{geometry}
\usepackage[T1]{fontenc}
\usepackage{lmodern}
\usepackage{amsmath,amssymb,amsthm}
\usepackage{microtype}
\usepackage{booktabs}
\usepackage[hidelinks]{hyperref}
\hypersetup{pdftitle={Exact Dyadic Tensor Products for the Newman Source},
            pdfauthor={Edward Baker}}

\newtheorem{theorem}{Theorem}[section]
\newtheorem{lemma}[theorem]{Lemma}
\newtheorem{proposition}[theorem]{Proposition}
\DeclareMathOperator{\sinc}{sinc}
\newcommand{\R}{\mathbb R}
\newcommand{\dd}{\mathop{}\!d}
\newcommand{\ee}{\mathrm e}
\newcommand{\ii}{\mathrm i}
\newcommand{\abs}[1]{\lvert#1\rvert}
\newcommand{\norm}[1]{\lVert#1\rVert}

\title{Exact Dyadic Tensor Products for the Newman Source\\
\large A uniform strip approximation by exact cell integration}
\author{Edward Baker}
\date{14 November 2026}

\begin{document}
\maketitle

\begin{abstract}
We represent the complete even theta source for the Newman field by
dyadic refinement. Conditional cell masses give isometries between
levels, and the corresponding square-root densities converge in a
weighted Hilbert space. Integrating the unresolved Fourier phase on
each cell introduces an exact sinc factor. A weighted derivative bound
for the complete source gives an absolute field error below
\(5\cdot10^{-7}\) with cutoff \(U=2\) and level \(L=30\), uniformly for
all real heights, imaginary parts of magnitude at most \(8\), and
\(0\le\lambda\le1/2\). At each fixed level the matching center-point
lattice has a uniform real-axis error at least half the total source
mass. The strip estimate is for exact cell masses and exact readout.
\end{abstract}

\section{Introduction}

At dyadic spacing \(h\), replace the positive Newman source by its
mean on each cell and retain the full Fourier phase inside the cell.
The weighted \(L^1\) error of the source then controls the field
even where \(|x|h\) is large. Binary scale indices are standard in
finite quantics representations \cite[\S II.1, (1)]{Shinaoka}, and
can encode polynomial-exact interpolation on finite grids
\cite[\S 4.1, Proposition 13]{Lindsey}. Here exact cell masses define
diagonal refinement factors and isometries between successive spaces
of square-root amplitudes.

\section{The complete theta source}\label{sec:source}

For \(t>0\) set \(\vartheta(t)=\sum_{n\in\mathbb Z}\ee^{-\pi n^2t}\).
Write
\[
 p(u)=\frac{\ee^u}{2}\bigl(\vartheta(\ee^{4u})-1\bigr),\qquad
 r_n(u)=\pi n^2\ee^{4u},\qquad
 \Phi(u)=\frac{p''(u)-p(u)}8.
\]
The series for \(p\) and every fixed derivative converge locally
uniformly on \(\R\). On a compact interval the derivative summands
are bounded by a polynomial in \(n\) times \(\ee^{-c n^2}\) with
\(c>0\), by differentiating \(\ee^u\ee^{-r_n(u)}\) and using
\(r_n'=4r_n\). Jacobi's transformation
\(\vartheta(t)=t^{-1/2}\vartheta(1/t)\)
\cite[20.2.3 and 20.7.32]{DLMFTheta} gives
\begin{equation}\label{eq:reflection}
 p(-u)-p(u)=\sinh u,\qquad p'(0)=-\tfrac12 .
\end{equation}
Indeed, \(p(-u)=\frac12(\ee^u\vartheta(\ee^{4u})-\ee^{-u})\);
differentiating the difference at zero gives the endpoint derivative.
Applying \(\partial_u^2-1\) to \eqref{eq:reflection} makes \(\Phi\)
even. Differentiating the series for \(p\) also gives
\begin{equation}\label{eq:theta-source}
 \Phi(u)=\ee^u\sum_{n\ge1}\bigl(2r_n(u)^2-3r_n(u)\bigr)\ee^{-r_n(u)}.
\end{equation}
Every summand is positive for \(u\ge0\) because \(r_n(u)\ge\pi>3/2\).
Evenness gives strict positivity on \(\R\). The reflection was derived
for the complete theta sum and ensures smoothness at zero.

Put
\begin{equation}\label{eq:conventions}
 W_\lambda(u)=\ee^{\lambda u^2}\Phi(u),\qquad
 H_\lambda(z)=\frac12\int_{\R} W_\lambda(u)\ee^{\ii zu}\dd u,\qquad
 \tau=\frac12-\lambda .
\end{equation}
We use \(0\le\lambda\le1/2\) unless specified otherwise, and write
\(z=x+\ii y\). The unnormalized mass
\(Z_\lambda=\int_{\R}W_\lambda\dd u=2H_\lambda(0)\) is strictly
positive. By evenness the readout is also the half-line cosine integral.

\begin{lemma}[Complete source and arithmetic tails]\label{lem:source}
For \(d,N\) nonnegative integers define
\[
 C_d(N)=\sum_{n>N}n^{2d}\ee^{-\pi(n^2-1)}.
\]
Let \(\rho=\pi\ee^{4|u|}\), \(Q_0(r)=2r^2-3r\), and
\(Q_{j+1}(r)=(1-4r)Q_j(r)+4rQ'_j(r)\).
If \(a_j\) is the sum of the absolute values of the coefficients of
\(Q_j\) and \(B_j=a_j C_{j+2}(0)\), then
\begin{equation}\label{eq:complete-dom}
 \abs{\Phi^{(j)}(u)}\le B_j\ee^{|u|}\rho^{j+2}\ee^{-\rho}.
\end{equation}
On \(u\ge0\) the same bound for the differentiated tail \(n>N\)
holds with \(B_j\) replaced by \(a_j C_{j+2}(N)\).
For every fixed \(b\ge0\) and every fixed nonnegative integer \(q\),
\[
 \sup_{0\le\lambda\le1/2}\int_{\R}
   \ee^{b|u|}(1+|u|^q)\bigl(W_\lambda(u)+|W'_\lambda(u)|\bigr)\dd u
   <\infty .
\]
\end{lemma}

\begin{proof}
The differentiated summand in \eqref{eq:theta-source} is
\(\ee^u Q_j(n^2\rho)\ee^{-n^2\rho}\) on the positive half-line.
The recurrence gives \(\deg Q_j\le j+2\), so
\[
 \abs{Q_j(n^2\rho)}\le a_j n^{2j+4}\rho^{j+2},\qquad
 \ee^{-n^2\rho}\le\ee^{-\rho}\ee^{-\pi(n^2-1)}.
\]
These inequalities prove the tail and full half-line bounds.
Evenness of \(\Phi\) gives \eqref{eq:complete-dom} everywhere.
Each \(C_d(N)\) is finite and tends to zero as \(N\to\infty\),
since successive terms eventually have ratio at most \(1/2\).
For the integral assertion, use \eqref{eq:complete-dom} with
\(j=0,1\) and
\[
 |W_\lambda'|\le \ee^{u^2/2}\bigl(|\Phi'|+|u|\Phi\bigr).
\]
On \(u\ge0\), substitution \(\rho=\pi\ee^{4u}\) gives
\(u=\tfrac14\log(\rho/\pi)\), \(\dd u=\dd\rho/(4\rho)\).
Apart from powers of \(\rho\) and \(\log(\rho/\pi)\), the integrand
contains \(\exp(\tfrac1{32}\log^2(\rho/\pi)-\rho)\), which is
bounded by \(\ee^{-\rho/2}\) for sufficiently large \(\rho\).
Evenness of the absolute terms gives the other half-line.
The argument also works with \(u^2/2\) replaced by \(D u^2\) for
any finite \(D\ge0\).
\end{proof}

These estimates imply joint entireness of the integral in
\((\lambda,z)\) and justify each fixed derivative under the integral
on compact parameter sets. In particular
\(\partial_\lambda H_\lambda=-\partial_z^2H_\lambda\), and in reverse
time \(\partial_\tau W_{1/2-\tau}=-u^2W_{1/2-\tau}\).

The normalization in \eqref{eq:conventions} can be seen directly.
Set \(J(z)=\int_0^\infty p(u)\cos(zu)\dd u\). The positive
half-line formula for \(p\) gives double-exponential decay of
\(p\) and \(p'\), by the same estimates on the summands as in
Lemma~\ref{lem:source}. Integrating \(p''\) twice by parts and
using \eqref{eq:reflection} yields
\begin{equation}\label{eq:endpoint}
 8H_0(z)=\tfrac12-(1+z^2)J(z).
\end{equation}
For \(\Re s>1\), termwise integration of the theta series gives
\(\pi^{-s/2}\Gamma(s/2)\zeta(s)
 =\int_0^\infty t^{s/2-1}(\vartheta(t)-1)/2\dd t\).
Indeed, the series of integrals is absolutely convergent for
\(\Re s>1\) by the gamma integral and \(\sum_{n\ge1}n^{-\Re s}<\infty\).
Split this integral at \(1\) and write
\[
 \frac{\vartheta(t)-1}{2}
   =t^{-1/2}\frac{\vartheta(1/t)-1}{2}+\frac{t^{-1/2}-1}{2}
 \qquad (0<t<1).
\]
The reflected term, after \(t\mapsto 1/t\), adds the power
\(t^{(1-s)/2-1}\) to the integral over \([1,\infty)\).
The elementary term integrates to \(1/(s-1)-1/s\).
Multiplying by \(s(s-1)/2\) thus gives
\[
 \xi(s)=\frac12+\frac{s(s-1)}2\int_1^\infty
  \frac{\vartheta(t)-1}{2}\bigl(t^{s/2-1}+t^{(1-s)/2-1}\bigr)\dd t
\]
in the convention of \cite[25.4.4]{DLMFZeta}. The expression extends
to all \(s\) by entireness of the integral. At \(s=(1-\ii z)/2\),
the substitution \(t=\ee^{4u}\) gives \(8J(z)\) for that integral,
and \(4s(s-1)=-(1+z^2)\). Comparing with \eqref{eq:endpoint} shows
\begin{equation}\label{eq:xi}
 H_0(z)=\tfrac18\,\xi(\tfrac12-\tfrac{\ii z}2).
\end{equation}

\section{Mass refinement and the square-root state}\label{sec:refinement}

For an integer \(L\ge0\) set \(h=2^{-L}\) and
\(I_{L,j}=[jh,(j+1)h)\), \(j\in\mathbb Z\), with center \(c_{L,j}\).
Write \(P_h f\) for the average of \(f\) on each such cell.
Let
\[
 m_{L,j}=\int_{I_{L,j}} W_\lambda(u)\dd u,\qquad
 f_L=P_h W_\lambda .
\]
The \(\lambda\)-dependence of \(m_{L,j}\) and \(f_L\) is suppressed.
The two children \(I_{L+1,2j}\), \(I_{L+1,2j+1}\) partition
\(I_{L,j}\). Their masses sum to \(m_{L,j}\), and all cell masses
are strictly positive.

There is a natural space for the common limit. On \(L^2(\R,\dd u)\),
let \(E_L\) consist of functions constant on every level-\(L\) cell.
The isometry from a cell coefficient vector \(v\) onto \(E_L\) is
\[
 T_L v=\sum_{j\in\mathbb Z}v_j h^{-1/2}\mathbf 1_{I_{L,j}} .
\]
Group \(j\) in the form \(j=2^L a+r\), with \(a\in\mathbb Z\) and
\(0\le r<2^L\). Writing \(r\) in \(L\) binary digits identifies the
coefficient space with
\(\mathcal H_L=\ell^2(\mathbb Z)\otimes(\mathbb C^2)^{\otimes L}\).
The first coordinate labels the unit interval. Appending the
reference bit \(\lvert+\rangle=(\lvert0\rangle+\lvert1\rangle)/\sqrt2\)
means that each basis vector \(e_j^{(L)}\) maps to
\((e_{2j}^{(L+1)}+e_{2j+1}^{(L+1)})/\sqrt2\).
Under the maps \(T_L\), this embedding is the inclusion
\(E_L\subset E_{L+1}\). The Hilbert-space completion of the nested
spaces is \(L^2(\R,\dd u)\). For instance, cell averages of a
continuous function of compact support approach it uniformly, stay
in a fixed bounded interval as \(h\to0\), and hence converge in
\(L^2\). Such functions are dense. Equivalently, the completion
uses the binary tensor product formed with \(\lvert+\rangle\)
as the reference vector at all unresolved bits.

The conditional mass isometry is different from the embedding used
to define this common space. At fixed \(\lambda\) it acts by
\begin{equation}\label{eq:mass-isometry}
 S_L e_j^{(L)}=\sum_{\epsilon=0}^1
       \sqrt{\frac{m_{L+1,2j+\epsilon}}{m_{L,j}}}
       e_{2j+\epsilon}^{(L+1)} .
\end{equation}
The two coefficients have squared sum one, and images of different
parents have disjoint support, so \(S_L\) extends to an isometry.
If \(v_{L,j}=\sqrt{m_{L,j}}\), then \(S_L v_L=v_{L+1}\) exactly.
Each vector has squared norm \(Z_\lambda\), and
\(T_L v_L=\sqrt{f_L}\). In the common \(L^2\) space the two
embeddings produce different functions: \(S_L\) updates the
density using the source masses, while appending
\(\lvert+\rangle\) leaves the old step function in place.

The corresponding diagonal products are explicit. On a fixed
unit interval \([a,a+1)\) with integer \(a\), let
\(R_k=f_k/f_{k-1}\) for \(k\ge1\). The ratio on a child is twice
its mass fraction in the parent. Thus \eqref{eq:mass-isometry}
can be applied by appending \(\lvert+\rangle\) and then
multiplying by \(\sqrt{R_{L+1}}\). The gate \(R_k\) acts on the
first \(k\) bits of the interval. The products telescope:
\begin{equation}\label{eq:telescopes}
  f_L=f_0\prod_{k=1}^L R_k,\qquad
  \sqrt{f_L}=\sqrt{f_0}\prod_{k=1}^L\sqrt{R_k}.
\end{equation}
For this interval, let \(m>0\) be the minimum of \(W_\lambda\) on
\([a,a+1]\times[0,1/2]\) and \(M<\infty\) the maximum of
\(\abs{W'_\lambda}\) there. The child and parent averages differ
by at most the parent oscillation, bounded by \(M2^{1-k}\).
It follows that
\[
 \norm{R_k-1}_{L^\infty([a,a+1])}\le \frac{M}{m}2^{1-k},
 \qquad
 \norm{\sqrt{R_k}-1}_{L^\infty([a,a+1])}
       \le \norm{R_k-1}_{L^\infty([a,a+1])}.
\]
Also \(\norm{f_L-W_\lambda}_{L^\infty([a,a+1])}\le M h\),
interpreting \(L^\infty\) in the essential sense at endpoints.
These estimates give local operator-norm convergence of the
multiplication products in \eqref{eq:telescopes}, uniformly in
\(\lambda\) on the specified time interval. The constants \(m,M\)
are local to each interval.

For the whole line define
\[
 \norm{g}_{1,b}=\int_{\R}\ee^{b|u|}|g(u)|\dd u,\qquad
 \norm{g}_{2,b}^2=\int_{\R}\ee^{b|u|}|g(u)|^2\dd u\qquad (b\ge0).
\]
Weighted convergence of \(\sqrt{f_L}\) follows from the following
estimate.

\begin{lemma}[Weighted projection]\label{lem:projection}
If \(f\) is locally absolutely continuous with
\(\norm{f}_{1,b}+\norm{f'}_{1,b}<\infty\), then
\begin{equation}\label{eq:projection}
 \norm{f-P_h f}_{1,b}\le \frac{h\ee^{bh}}2\norm{f'}_{1,b},
 \qquad
 \norm{P_h f}_{1,b}\le \ee^{bh}\norm{f}_{1,b}.
\end{equation}
\end{lemma}

\begin{proof}
On \(I=[a,a+h]\) with mean \(\bar f\), the triangle inequality
and absolute continuity imply
\[
 \begin{split}
 \int_I|f(v)-\bar f|\dd v
 &\le h^{-1}\int_I\int_I |f(v)-f(w)|\dd w\dd v\\
 &\le \frac{2}{h}\int_a^{a+h}(t-a)(a+h-t)|f'(t)|\dd t .
 \end{split}
\]
The two factors in the kernel give the measure of the choices
of one endpoint on each side of a fixed \(t\). The factor of two
counts both orders. Since this kernel is at most \(h/2\), the
unweighted estimate follows on \(I\). The ratio of maximum to
minimum of \(\ee^{b|u|}\) on \(I\) is at most \(\ee^{bh}\).
Multiplying the unweighted estimate by the maximum and comparing
it to the minimum proves the first inequality after summing.
For the second, use \(h|\bar f|\le\int_I|f|\) on each cell and
the same comparison of weights.
\end{proof}

For \(q,r\ge0\), \(|\sqrt q-\sqrt r|^2\le |q-r|\).
Lemma~\ref{lem:source} and Lemma~\ref{lem:projection} give
\begin{equation}\label{eq:state-limit}
 \norm{T_L v_L-\sqrt{W_\lambda}}_{2,b}^{\,2}
 \le \norm{f_L-W_\lambda}_{1,b}
 \le \frac{h\ee^{bh}}2\norm{W_\lambda'}_{1,b}.
\end{equation}
Thus the prepared state converges in every fixed weighted
\(L^2\) space, uniformly in \(\lambda\). The result with \(b=0\)
identifies the limit in the reference-bit tensor product, and
the weighted norm further controls its tail. These are statements
about the prepared states and densities on the complete line.
The multiplication-product norm estimate above is interval by
interval.

\section{Phase on the binary coordinates}\label{sec:readout}

On \([a,a+1)\) write \(u=a+\sum_{k\ge1}2^{-k}b_k\) with
\(b_k\in\{0,1\}\) almost everywhere. With \(n_k\) denoting the
projector onto bit \(1\) in coordinate \(k\), the multiplier
\(\ee^{\ii zu}\) on this interval has the product
\begin{equation}\label{eq:phase-product}
 \ee^{\ii za}\prod_{k\ge1}\exp(\ii z 2^{-k}n_k).
\end{equation}
For \(|z|\le D\) the \(k\)-th bit factor differs from the identity
in norm by at most \(D2^{-k}\ee^{D2^{-k}}\), using
\(\ee^w-1=w\int_0^1\ee^{tw}\dd t\). Partial products act by
the corresponding truncated exponential, whose coordinate differs
from \(u\) by at most \(2^{-k}\) after retaining \(k\) bits.
It follows directly that the products converge in local
multiplication-operator norm, uniformly on compact \(z\)-sets
for a fixed interval. For any \(k\), a real choice \(x=\pi2^k\)
makes the \(k\)-th factor differ from the identity by \(2\).
Uniform convergence to the identity fails over real \(x\).

On \(I_{L,j}\) write \(u=jh+s\). Under normalized Lebesgue
measure on the cell, the unresolved bits encode \(s\) in
\([0,h)\) and have uniform binary-prefix probabilities: specifying
\(r\) new bits restricts \(s\) to an interval of length \(h2^{-r}\).
Their product phase therefore has conditional mean
\begin{equation}\label{eq:sinc}
 \frac1h\int_0^h \ee^{\ii zs}\dd s
   =\ee^{\ii zh/2}\sinc(zh/2),\qquad
 \sinc(w)=\begin{cases}\sin(w)/w,&w\ne0,\\1,&w=0.\end{cases}
\end{equation}
The limit under the mean is justified by the compact-\(z\) convergence
in \eqref{eq:phase-product}. Equation \eqref{eq:sinc} is an exact
identity for every \(z\).

Define the readout
\begin{equation}\label{eq:readout}
 \begin{split}
 A_L(\lambda,z)
  &=\frac12\int_{\R} f_L(u)\ee^{\ii zu}\dd u\\
  &=\frac12\sinc(zh/2)\sum_{j\in\mathbb Z}
                   m_{L,j}(\lambda)\,\ee^{\ii zc_{L,j}}.
 \end{split}
\end{equation}
This is also half the matrix element of the phase in the
prepared state \(\sqrt{f_L}\). For complex \(z\) this state is
in the domain of the phase multiplier, by the second inequality
of \eqref{eq:projection} with \(b=2|\Im z|\).
Each term of the integral over a single cell gives
\eqref{eq:sinc} times its mass and left-endpoint phase.
For \(|y|\le b\), summing is justified by
\[
 \sum_j m_{L,j}\,\abs{\ee^{\ii zc_{L,j}}}
 \le \ee^{bh/2}\norm{W_{1/2}}_{1,b}.
\]
Indeed \(|c_{L,j}|\le |u|+h/2\) throughout the cell.
The same estimate for the tails of the series proves uniform
convergence on compact sets of \(z\), so \(A_L\) is entire.
For an integer \(U\ge1\) let \(A_{L,U}\) use just the cells
inside \([-U,U]\), with indices \(-U/h\le j<U/h\).
The same integral and contraction descriptions hold with that
restriction.

\section{Uniform absolute strip error}\label{sec:strip}

\begin{theorem}\label{thm:strip}
Define
\[
 V_8=\sup_{0\le\lambda\le1/2}\norm{W_\lambda'}_{1,8},
 \qquad
 T_8(U)=\frac12\int_{|u|>U}\ee^{8|u|}W_{1/2}(u)\dd u .
\]
Then \(V_8<868\). For all \(0\le\lambda\le1/2\), \(x\in\R\),
and \(|y|\le8\), one has
\begin{align}
 |H_\lambda(x+\ii y)-A_L(\lambda,x+\ii y)|
     &\le \frac{h\ee^{8h}}4 V_8, \label{eq:full-bound}\\
 |H_\lambda(x+\ii y)-A_{L,U}(\lambda,x+\ii y)|
     &\le \frac{h\ee^{8h}}4 V_8+T_8(U). \label{eq:cut-bound}
\end{align}
For \(U=2\) and \(L=30\), the right side of
\eqref{eq:cut-bound} is less than \(5\cdot10^{-7}\).
\end{theorem}

\begin{proof}
For \eqref{eq:full-bound}, apply \eqref{eq:projection}
to \(W_\lambda\) with \(b=8\), and then use the Fourier factor
\(1/2\) and \(|\ee^{\ii zu}|\le \ee^{8|u|}\) in the difference
of the two integrals. At the cutoff apply the same cell estimates
on \([-U,U]\). The exterior of the true integral is bounded by
\(T_8(U)\), since \(0<W_\lambda\le W_{1/2}\).
The estimate uses the source variation but no bound involving
the real part \(x\).

For the explicit constants work first on \(u\ge0\).
The recurrence in Lemma~\ref{lem:source} gives
\[
 Q_1(r)=-8r^3+30r^2-15r,\qquad
 B_0=5C_2(0),\qquad B_1=53C_3(0).
\]
For \(a_n=n^6\ee^{-\pi(n^2-1)}\) and \(n\ge2\),
\[
 \frac{a_{n+1}}{a_n}
 =\left(\frac{n+1}{n}\right)^6\ee^{-\pi(2n+1)}
 \le (3/2)^6\ee^{-15}<\tfrac12,\qquad
 a_2=64\ee^{-3\pi}<\tfrac18.
\]
Here \((3/2)^6<12\), \(\pi\ge3\), and \(\ee>2\).
Thus \(C_2(0)\le C_3(0)\le 1+2a_2<2\).
Consequently \(B_0<10\), \(B_1<106\). By
\eqref{eq:complete-dom} and \(2\lambda\le1\),
\begin{equation}\label{eq:derivative-source}
 |W_\lambda'(u)|\le
    \ee^{u^2/2+u}(B_1\rho^3+u B_0\rho^2)\ee^{-\rho}.
\end{equation}
The exponential series gives \(\rho\ge 3(1+4u+8u^2)\) and
\(\rho\ge 3(1+4u)\). In particular,
\[
 \frac{\rho}{2}-\frac{u^2}{2}-9u
 \ \ge\ \frac{23}{2}u^2-3u+\frac32
 \ =\ \frac{23}{2}(u-3/23)^2+\frac{30}{23}>0,\qquad
 u\le \frac{\rho}{4}.
\]
Multiplying \eqref{eq:derivative-source} by \(\ee^{8u}\)
gives \((B_1+B_0/4)\rho^3\ee^{-\rho/2}\) as an upper bound.
Evenness of \(|W_\lambda'|\) and the change of variable
\(\dd u=\dd\rho/(4\rho)\) yield
\[
 V_8\le \frac{B_1+B_0/4}{2}
             \int_\pi^\infty \rho^2\ee^{-\rho/2}\dd\rho
     \le 8B_1+2B_0<868.
\]
In the second inequality the integral was extended to zero,
where its value is \(16\), by the substitution \(\rho=2t\)
and two integrations by parts.

Write \(\rho_U=\pi\ee^{4U}\).
The Fourier factor cancels the doubling from the two half-lines.
Using the \(j=0\) source bound and the same exponential domination
on \(u\ge U\) gives
\begin{equation}\label{eq:tail}
 T_8(U)\le \frac{B_0}{4}\int_{\rho_U}^\infty
             \rho\,\ee^{-\rho/2}\dd\rho
          =\frac{B_0}{2}(\rho_U+2)\ee^{-\rho_U/2}.
\end{equation}
For \(U=2\), the exponential series gives
\(\rho_2\ge3(1+8+8^2/2)>64\). The function
\((r+2)\ee^{-r/2}\) decreases for \(r>0\), so
\(T_8(2)<5(64+2)\ee^{-32}<330/2^{32}\).
For \(h=2^{-30}\), \(8h<1/2\) and
\(\ee^{8h}\le(1-8h)^{-1}<2\), comparing the exponential series
to the geometric series. Therefore the cell term in
\eqref{eq:cut-bound} is less than \(1736/2^{32}\). Combining
the bounds,
\[
 \frac{h\ee^{8h}}4 V_8+T_8(2)
   <\frac{2066}{2^{32}}<\frac{1}{2\cdot10^6},
\]
where the last comparison is
\(4{,}132{,}000{,}000<4{,}294{,}967{,}296\).
\end{proof}

\section{Aliasing of the fixed center-point lattice}\label{sec:alias}

An alternative scalar rule with the same masses evaluates the
phase only at the cell centers:
\[
 B_L(\lambda,x)=\frac12\sum_{j\in\mathbb Z}
                 m_{L,j}(\lambda)\,\ee^{\ii x c_{L,j}}
 \qquad (x\in\R).
\]

\begin{proposition}\label{prop:aliasing}
For every fixed \(L\) and \(0\le\lambda\le1/2\),
\[
 \sup_{x\in\R}|H_\lambda(x)-B_L(\lambda,x)|\ \ge\ Z_\lambda/2 .
\]
\end{proposition}

\begin{proof}
At \(x_k=2\pi k/h\), \(k\) a positive integer,
\(\ee^{\ii x_k c_{L,j}}=(-1)^k\) for every \(j\).
Hence \(B_L(\lambda,x_k)=(-1)^k Z_\lambda/2\) by absolute
summability. The complete source decays at infinity and its
derivative is integrable by Lemma~\ref{lem:source}. Integration
by parts gives, for \(x\ne0\),
\[
  |H_\lambda(x)|\le \norm{W_\lambda'}_{1,0}/(2|x|).
\]
Thus the absolute error at \(x_k\) tends to
\(Z_\lambda/2\). Taking the supremum over all real \(x\) proves
the bound.
\end{proof}

The proposition concerns the supremum over the full real axis
for the fixed-level center-point rule with full-line masses.
At the indicated frequencies \(\sinc(x_kh/2)=0\), which exactly
cancels the lattice phase in \(A_L\). Directly integrating the
same step density gives the identical field \eqref{eq:readout},
including the same cutoff when applied. This is the scalar
baseline for the exact-cell construction.

\section{Time and arithmetic in the finite readout}\label{sec:time}

The exact masses at each time are integrals of
\(W_{1/2-\tau}=\ee^{-\tau u^2} W_{1/2}\). Differentiation at fixed
cell and time gives
\[
 \partial_\tau m_{L,j}(1/2-\tau)
    =-\int_{I_{L,j}}u^2 W_{1/2-\tau}(u)\dd u .
\]
After dividing by \(h\) this yields
\(\partial_\tau P_h W_{1/2-\tau}=-P_h(u^2 W_{1/2-\tau})\).
The latter is constant on each cell, whereas
\(-u^2 P_h W_{1/2-\tau}\) varies inside every positive-mass cell.
Thus the finite refinement follows the exact orbit by projection
at each time. Its time derivative uses further source information
within the cells.

Theorem~\ref{thm:strip} assumes the complete masses and exact
contraction. If the retained masses instead have errors
\(\delta m_{L,j}(\lambda)\), the centered average representation
\[
 \sinc(zh/2)=h^{-1}\int_{-h/2}^{h/2}\ee^{\ii zv}\dd v
\]
bounds its modulus by \(\ee^{4h}\) on \(|y|\le8\).
If \(e_{\rm readout}\) is an absolute bound for arithmetic in
evaluating \eqref{eq:readout} at the supplied masses, the total
computed-minus-ideal error is at most
\begin{equation}\label{eq:computed}
 \frac12\ee^{4h}\sum_{-U/h\le j<U/h}
             \ee^{8|c_{L,j}|}|\delta m_{L,j}(\lambda)|
       +e_{\rm readout}.
\end{equation}
Each mass error must include the contributions of any theta-series
truncation, quadrature, and approximation of the refinement factors
used to obtain that mass. For a computed error below \(10^{-6}\)
at the theorem's tuple, it suffices to bound the weighted sum in
\eqref{eq:computed} by \(10^{-6}/(2\ee^{4h})\) and bound
\(e_{\rm readout}\) by \(10^{-6}/4\). The first bound is needed
uniformly in \(\lambda\), and the second over the same time and
strip domain as the theorem. These are requirements for an
implementation. The exact source and contraction in the theorem
carry only the paid projection and exterior errors.

\section{Outstanding challenges}

The cutoff at the given resolution specifies \(2U/h=2^{32}\)
cell masses. Obtaining a compressed description of those complete
masses, with a proved error in the weighted sum
\eqref{eq:computed}, remains open. A useful cost result would
also bound the ranks of the growing-prefix refinement factors
and the cost of their contraction, and compare those costs with
direct integration of the same piecewise-constant density.
An arithmetic theorem for unbounded real \(x\) would need
phase evaluation with a justified uniform error over that domain.

For the full-line product the convergence established here is
weighted state convergence. An operator-norm theorem for the
full-line source factors would require new control of the
relative child and parent averages far out in the theta tail.
The local bounds based on a positive interval minimum leave
that question open. For time stepping, a finite closed evolution
would require information about the moments inside each cell or
an estimate replacing that information with a controlled error.

Finally, the strip estimate is absolute in the physical
normalization \eqref{eq:conventions}. A relative estimate near
zeros or a signed consequence for a logarithmic Green field
would require additional lower or deflation estimates. Weighted
regularity and the exact readout apply to smooth sources that
admit collisions. The quantitative theta input used here is
the complete source bound for its variation and tail.

\paragraph{Acknowledgment.} This draft was prepared for Edward Baker with substantial
OpenAI Codex LLM assistance. Recorded worker dispatch was GPT-6.1 Sol /
Ultra. Independent backend metadata are unavailable.

\begin{thebibliography}{9}

\bibitem{DLMFTheta}
W.~P. Reinhardt and P.~L. Walker,
\emph{Theta functions}, NIST Digital Library of Mathematical Functions,
Chapter 20, \href{https://dlmf.nist.gov/20.2#E3}{Eq.~20.2.3} and
\href{https://dlmf.nist.gov/20.7#E32}{Eq.~20.7.32}.

\bibitem{DLMFZeta}
T.~M. Apostol, \emph{Zeta and related functions},
NIST Digital Library of Mathematical Functions, Chapter 25,
\href{https://dlmf.nist.gov/25.4#E4}{Eq.~25.4.4}.

\bibitem{Lindsey}
M.~Lindsey, \emph{Multiscale interpolative construction of quantized
tensor trains}, \href{https://arxiv.org/abs/2311.12554v3}{arXiv:2311.12554v3},
2024.

\bibitem{Shinaoka}
H.~Shinaoka, M.~Wallerberger, Y.~Murakami, K.~Nogaki, R.~Sakurai,
P.~Werner, and A.~Kauch, \emph{Multiscale space-time ansatz for
correlation functions of quantum systems based on quantics tensor
trains}, \href{https://doi.org/10.1103/PhysRevX.13.021015}
{Phys.\ Rev.\ X \textbf{13}, 021015 (2023)}.

\end{thebibliography}

\end{document}
